% =====================================================================
% Label Semantic Aware Pre-training for Few-shot Text Classification
% Reconstructed, editable LaTeX source.
%
% NOTE ON PROVENANCE
%   The repository originally contained only the compiled paper.pdf
%   (built with pdfTeX / "LaTeX with hyperref", 2023-05-20); the LaTeX
%   source was never committed. This file reconstructs the paper from the
%   text of that PDF and applies a set of corrections to the prose.
%
%   Every substantive change to the original wording is tagged with a
%   "% CORRECTED:" comment so the authors can audit each edit. A summary
%   of all corrections is in paper/CORRECTIONS.md.
%
%   Figures 1-8 in the original are raster images that are NOT present in
%   the repository, so they are reproduced here as captioned placeholders
%   (see the \figplaceholder macro). Drop the original images into
%   paper/figures/ and replace the placeholders to restore them.
%
% Compiles with a plain pdflatex run (article class; no ACL style files
% required) so the authors can edit and rebuild without external deps.
% =====================================================================
\documentclass[11pt]{article}

\usepackage[T1]{fontenc}
\usepackage[utf8]{inputenc}
\usepackage[margin=1in]{geometry}
\usepackage{graphicx}
\usepackage{booktabs}
\usepackage{amsmath}
\usepackage{authblk}
\usepackage{url}
\usepackage{xcolor}
\usepackage[hidelinks]{hyperref}

% Placeholder for figures whose source images are not in the repo.
\newcommand{\figplaceholder}[1]{%
  \begin{center}
  \fbox{\parbox{0.8\linewidth}{\centering\vspace{1.5em}%
  \textit{[Figure source image not available in repository.]}\\[0.3em]%
  \footnotesize #1\vspace{1.5em}}}
  \end{center}}

\title{Label Semantic Aware Pre-training for Few-shot Text Classification}

\author[ ]{Veer Singh}
\author[ ]{Venkata Sai Phanindra}
\author[ ]{Lokesh Tangella}
\author[ ]{Sanjana Radhakrishna}
\affil[ ]{University of Massachusetts Amherst\\
\texttt{\{veersingh, vpamidimukkala, ltangella, sanjanaradha\}@umass.edu}}

\date{}

\begin{document}
\maketitle

% CORRECTED: Added an explicit replication-study framing up front so the
% paper does not later mis-describe LSAP as the authors' own novel method.
\begin{abstract}
This is a \emph{replication and analysis study} of the Label Semantic Aware
Pre-training (LSAP) technique introduced by Mueller and Roth (2022). We
re-implement LSAP on \mbox{T5-small} (60M parameters) rather than the
multi-billion-parameter models used in the original work, and evaluate
few-shot intent classification across 1-, 2-, 4-, 8-, 16-shot, and
full-resource settings on ATIS, SNIPS, and TOPv2. We confirm the original
paper's qualitative trend---label denoising is the most effective secondary
pre-training format and helps most in the lowest-resource settings---while
also reporting where our small-scale reproduction \emph{diverges} from the
original: the few-shot accuracy advantage of the LSAP model is smaller and
noisier at this scale, plain T5-small remains highly competitive (and
superior on several generation-quality metrics), and from-scratch
pre-training on our 120K-example corpus failed to produce usable models.
\end{abstract}

\section{Introduction}

The advancement of large pre-trained language models has greatly enhanced
performance on natural language processing (NLP) tasks, particularly in
scenarios with limited training data. However, while these models excel at
encoding input information, there has been limited investigation into
providing informative representations of \emph{labels} to the models. Most
text classification methods commonly utilize label indices, but research
suggests that employing sequence-to-sequence models to generate labels can
improve performance. Nonetheless, these generative models only leverage
label semantics during fine-tuning and prediction, rather than during
pre-training.

To address this gap, a recent study introduced the Label Semantic Aware
Pre-training (LSAP) approach (Mueller and Roth, 2022), which integrates
label semantics and input--label associations into the pre-training phase.
This approach has shown the potential to achieve superior performance with a
reduced number of fine-tuning examples across diverse domains. Our research
aims to assess the effectiveness of LSAP by pre-training two generative
models, one of which incorporates label semantics and one of which does not.
We evaluate their few-shot performance on intent and topic classification
datasets. Furthermore, time permitting, we planned to extend this approach
to encoder models like BERT-base and evaluate their performance on sentiment
analysis tasks.

\section{Proposed \& Accomplished Goals}

Our principal aim is to implement the LSAP technique on the T5-small model
and evaluate its performance across diverse few-shot settings, comparing it
to our baseline models. We also incorporated a stretch goal alongside our
initial proposal. All of the planned core experiments were successfully
conducted, and the code can be found on
GitHub.\footnote{\url{https://github.com/DigitalVeer/nlp-project}} While we
made strides toward the stretch goal, it remains incomplete.

In the course of this investigation, we undertook a variety of tasks to
examine the effectiveness of different pre-training approaches and
fine-tuning strategies on downstream performance:

\begin{itemize}
  \item We gathered and preprocessed appropriate datasets for our study
        (Section~\ref{sec:data}).
  \item We explored three distinct pre-training formats, scrutinizing their
        performance under various shot settings (1, 2, 4, 8, 16-shot, and a
        full-resource scenario) for downstream tasks.
  \item We implemented the LSAP technique and evaluated its efficiency
        through secondary pre-training applied to a quartet of T5-small
        models.
  \item We performed fine-tuning with respect to hyperparameters for the
        downstream task of intent classification.
  \item We evaluated the primary pre-trained T5-small model and the four
        secondarily pre-trained models under several few-shot scenarios,
        from 1-shot to full-resource.
  \item Errors committed by the best-performing model in $k$-shot scenarios
        were scrutinized to gain a deeper understanding of the model's
        weaknesses.
  \item We made a comparative analysis of different metrics across all five
        models in the $k$-shot settings.
\end{itemize}

Initially, our study had a broader scope that included an ambitious stretch
goal. However, we encountered limitations in time and computational
resources that prevented us from achieving it. We operated on Google
Colaboratory, which offers limited GPU availability, constraining our
computational capabilities. The unattained stretch goal involved an in-depth
analysis of the effectiveness of the LSAP technique when applied to
encoder-only models, like BERT, for the downstream task of sentiment
analysis.

\section{Related Work}

Label semantics have been widely used to enhance performance and robustness
across settings and tasks, even prior to the prevalence of dense embedding
representations. For instance, Chang et al.\ (2008) achieved over 80\%
accuracy in binary text classification tasks without any labeled training
examples, by employing naive Bayes classifiers that utilized rich semantic
representations of labels (Gabrilovich and Markovitch, 2007). Similarly,
Song and Roth (2014) adopted a comparable strategy for hierarchical and
multinomial classification tasks and found that dataless approaches could
approach, and sometimes outperform, supervised approaches.

More recently, label semantics based on dense embeddings have gained
popularity, particularly with the advent of contextualized word embeddings
(Peters et al., 2018). Label-wise attention networks (LWANs) have been
proposed for datasets with large structured label spaces, where multiple
labels can apply to a long document. Mullenbach et al.\ (2018) and Rios and
Kavuluru utilized LWANs based on convolutional neural networks and extended
the attention mechanism for zero-shot settings. Chalkidis et al.\ (2020)
employed BERT to embed labels for an LWAN. In contrast, our setting involves
a smaller label space and shorter input texts. Another label-semantic-aware
approach, CLESS, performs contrastive pre-training on a question-answering
dataset with a large and sparse label space for massively multi-label text
classification. Our setting exhibits a larger discrepancy between
pre-training and fine-tuning, and we aim for more domain generalization by
leveraging token and span reconstruction objectives.

In short-text intent and topic classification, other works have integrated
label embeddings into their systems. Gaonkar et al.\ (2020) and Ma et al.\
(2022) used label embeddings from BERT and a label attention mechanism to
improve emotion classification accuracy and few-shot named entity
recognition, respectively. Generative approaches, such as those of Rongali
et al.\ (2020), Athiwaratkun et al.\ (2020), and Paolini et al.\ (2021),
implicitly leverage label semantics for text and token classification tasks
by generating labels at prediction time.

Few-shot text classification involves classifying text with only a few
examples from the training set. Recent approaches include
(Ro)BERT(a)-based methods (Chen et al., 2020; Yang et al., 2019),
prototypical networks, dynamic memory induction networks, and generative
classification. In our work, we evaluate the LSAP approach on two subtasks
of text classification: topic classification (TC) and intent classification
(IC). In IC, the input is a conversational utterance and the output is the
label representing the user's intended action. Recent few-shot IC approaches
include dual encoders (Cer et al., 2018; Casanueva et al., 2020), ConveRT
(Henderson et al., 2020), prototypical networks and meta-learning (Krone et
al., 2020), nearest-neighbor discriminative methods (Zhang et al., 2020),
and span-level contextualized embedding retrieval (Yu et al., 2021).

Generative text classification has become more feasible with recent
advancements in large pre-trained language models such as GPT (Radford et
al., 2019) and GPT-3 (Brown et al., 2020). This approach trains a model to
generate natural-language labels from an input sequence, reducing the
train--test gap without architectural changes. Pattern-exploiting training
(PET) formats training and testing examples as cloze-style prompts, where
the label is typically a single word. Our approach instead provides the
model with concatenated utterances and labels, allowing it to transduce
variable-length label sequences without reformatting the evaluation data.
Seq2seq approaches using pointer/copy mechanisms (Rongali et al., 2020) or
T5 (Raffel et al., 2020) have shown effectiveness and data efficiency in
text and token classification tasks. We build upon T5-small but, following
Mueller and Roth (2022), perform a label-semantic-aware secondary
pre-training step on a variety of datasets before fine-tuning.

\section{Data}
\label{sec:data}

% CORRECTED: corpus-size wording made internally consistent with Table 1
% (10,003 + ~110,000 ~= ~120,000 examples).
Our pre-training data consists of utterances with intent labels
(Table~\ref{tab:pretrain}), similar to the original LSAP paper. For
pre-training we use the same datasets as the original authors: the PolyAI
Banking dataset\footnote{\url{https://polyai.com/datasets/banking}} and the
WikiHow intents
dataset.\footnote{\url{https://github.com/zharry29/wikihow-intent}} The
PolyAI Banking dataset contains over 10,000 training examples categorized
into 77 unique intents, focusing on online banking queries.

\figplaceholder{Figure 1: Example utterance from PolyAI Banking.}

The WikiHow intents dataset consists of heuristically labeled examples
derived from WikiHow articles. Each example is a single line taken from an
article, typically the longest step, and its intent label is derived from
the article title by removing the phrase ``How To.''

\figplaceholder{Figure 2: Processing intents from WikiHow.}

The pre-training corpus combines these two datasets, resulting in a little
over 120,000 examples.

For evaluation, we used samples from
ATIS,\footnote{\url{https://github.com/howl-anderson/ATIS_dataset/}}
SNIPS,\footnote{\url{https://github.com/ZhenwenZhang/Slot_Filling}}
TOPv2 (reminder), and TOPv2 (weather). ATIS is commonly used for
benchmarking intent classification, and TOPv2 is common for task-oriented
semantic parsing. SNIPS and ATIS are balanced, while TOPv2 is unbalanced. We
chose these datasets because they are non-overlapping with our pre-training
corpus. All datasets used are publicly available.

We categorize the quality of the PolyAI Banking dataset as ``gold'' and the
WikiHow dataset as ``silver'' based on the labeling method. PolyAI Banking is
human-labeled (gold); WikiHow is labeled heuristically (silver).

\begin{table}[t]
\centering
\small
\begin{tabular}{lllrr}
\toprule
Dataset & Quality & Description & Train ex. & Intents \\
\midrule
PolyAI Banking & Gold & Online banking queries & 10,003 & 77 \\
WikiHow Intents & Silver & \begin{tabular}[t]{@{}l@{}}Article titles\\(``How To''\\removed)\end{tabular} & 110,000 & 66,343 \\
\bottomrule
\end{tabular}
\caption{Pre-training datasets. ``Quality'' refers to the source of the
labels (human-labeled is gold, deterministically labeled is silver).}
\label{tab:pretrain}
\end{table}

\subsection{Data preprocessing}

In the LSAP approach, a secondary pre-training step is performed using data
labeled with semantics. We consider/evaluate the following pre-training
formats:

\begin{enumerate}
  \item \textbf{Random Span Denoising:} Each intent is appended to the end
        of its respective utterance, then 15\% of the tokens in the
        utterance--intent input sequence are randomly masked; the goal is to
        reconstruct the noised spans in the output sequence. This is the
        same objective T5 uses.
  \item \textbf{Intent Classification:} We prefix ``intent classification:''
        to the start of every utterance and keep the output sequence as the
        intent, keeping the data in the same format used for supervised
        fine-tuning.
  \item \textbf{Label Denoising:} We always noise the entire intent label
        and append a mask to the end of each utterance; the pre-training
        objective is to reconstruct the output from the given mask. As the
        original authors note, this frames intent classification as an
        unsupervised denoising task.
\end{enumerate}

\figplaceholder{Figures 3--5: Random span denoising, intent classification,
and label denoising formats.}

As our results show, the choice of pre-training format has a large impact.
We implement different $n$-shot models across all three formats and concur
with the original authors that label denoising is the preferred format for
our downstream task.

\section{Baselines}

Our baselines consist of three models:
\begin{itemize}
  \item \textbf{Model 1:} T5-small with no secondary pre-training.
  \item \textbf{Model 2:} T5-small with secondary pre-training on utterances
        (no intents).
  \item \textbf{Model 3:} T5-small with secondary pre-training on utterances
        together with their intents (LSAP).
\end{itemize}

We evaluate these three models across multiple few-shot settings (1-, 2-,
4-, 8-, 16-shot, and full resource). A $k$-shot setting denotes fine-tuning
a pre-trained model using $k$ samples per label, including parameter
updates. The full-resource setting entails fine-tuning on the entirety of
the available labeled data.

Because Model 3 was pre-trained with intents, we anticipate favorable
adaptation to the downstream intent classification task even with few
labeled examples. We therefore expect Model 3 to outperform both T5-small
(Model 1) and Model 2 in low-shot settings; T5-small and Model 2 are
selected as our baselines. The expected ordering is
Model~1 $<$ Model~2 $<$ Model~3 in low-shot settings, with performance
converging as $k$ increases.

% CORRECTED: clearly mark Models 4/5 as a from-scratch ablation that, in our
% runs, failed (see Sections 6 and 6.1), so they are not over-sold here.
We additionally experimented with two T5-small-based models, Model 4 and
Model 5, both initialized with \emph{random weights} (erasing the
domain-specific knowledge obtained from pre-training) and then pre-trained
on our corpus following the same recipe as Models 2 and 3, respectively.
Model 5 (with intents) is expected to outperform Model 4 (without), so Model
4 is the baseline for Model 5. T5-small is still expected to outperform Model
5 because of its large-scale pre-training. As reported below, this
from-scratch ablation did not succeed at our scale.

During pre-training, the PolyAI Banking and WikiHow datasets were split into
training and validation sets in a 96:4 ratio (an 80:20 split would be ideal;
96:4 was chosen due to computational limits). The validation set was used to
compute metrics such as perplexity after each epoch. The fine-tuning
datasets followed a standard 90:5:5 train/test/validation split, with
stratification to ensure equal label proportions.

\subsection{Hyperparameter Tuning}

During secondary pre-training we adjusted the learning rate, weight decay,
number of epochs, training batch size, and evaluation batch size, following
the guidelines in the referenced paper. Due to resource limits we hit
out-of-memory issues and reduced the batch size from an initial 128 down to
as low as 1, which let the model run but greatly increased training time. We
then employed gradient accumulation (factor of 4) with an evaluation batch
size of 2 and a training batch size of 8 to complete epochs without memory
constraints. We logged perplexity and loss for both training and evaluation
throughout.

After further heuristic tuning we settled on: training batch size 16,
evaluation batch size 2, gradient accumulation steps 4, learning rate
$5\times10^{-3}$, and 5 training epochs. With these settings all models
reached train/validation perplexity in the 1--2 range, \emph{except} Models 4
and 5, which exhibited perplexity above 40 at the end of epoch 5. Extending
their training to 10 epochs reduced perplexity to roughly 20. The higher
perplexity of Models 4 and 5 is attributable to their random initialization,
which discards information from the large pre-training corpus.

In the fine-tuning / few-shot stage we tuned learning rate, weight decay,
number of epochs, training batch size, and evaluation batch size. Due to
limited compute and the time cost of each run, we adjusted hyperparameters
one at a time: holding all others fixed, observing the direction of change
in validation accuracy, and adjusting accordingly, repeating until results
were satisfactory. The test set was held out until all hyperparameters were
tuned. Typical ranges were:

\begin{enumerate}
  \item Learning rate: $[4\times10^{-3},\, 2\times10^{-2}]$.
  \item Batch size: $\{16, 32\}$.
  \item Weight decay: 0.
  \item Number of epochs: 3 or 10 (varied by dataset).
\end{enumerate}

% CORRECTED: flag the apparent inconsistency between the single LR used for
% secondary pre-training and the tuning range, so readers are not misled.
We note that the single learning rate above ($5\times10^{-3}$) refers to the
secondary pre-training stage, whereas the range applies to fine-tuning; both
are high for T5 and contributed to the run-to-run variance we observe.

\section{Approach}

% CORRECTED (key fix): the original text here read "we propose a novel
% approach called Label Semantic Aware Pre-training (LSAP)". LSAP is NOT our
% contribution -- it is from Mueller and Roth (2022). Reframed as a
% replication/analysis study to remove the misattribution.
Prior work has explored incorporating supplementary details, such as the
intent of input sequences, into models---but typically during the
fine-tuning phase. Comparatively little work uses label semantics during
\emph{pre-training}. This paper \emph{replicates and analyzes} the Label
Semantic Aware Pre-training (LSAP) approach of Mueller and Roth (2022),
which leverages label semantics during the pre-training phase to improve
few-shot performance. Our contribution is not the method itself but a
small-scale (T5-small) reproduction, an added suite of generation-quality
metrics, and an honest accounting of where the technique's benefits do and
do not transfer at this scale.

The goal is to analyze whether the few-shot performance of a language model
improves when intents are introduced during a secondary pre-training stage.
We consider five models, all based on T5-small (60M parameters):

\begin{enumerate}
  \item \textbf{Model 1 --- T5-small:} a pre-trained T5-small model.
  \item \textbf{Model 2 --- T5-small-Adapted:} T5-small with secondary
        pre-training on our corpus \emph{without} intents.
  \item \textbf{Model 3 --- T5-small-LSAP:} T5-small with secondary
        pre-training on our corpus \emph{with} intent labels.
  \item \textbf{Model 4 --- T5-small-Nomad-Adapted:} randomly initialized
        T5-small, then pre-trained on our corpus without intents.
  \item \textbf{Model 5 --- T5-small-Nomad-LSAP:} randomly initialized
        T5-small, then pre-trained on our corpus with intents.
\end{enumerate}

There are various ways to perform secondary pre-training (except for Model
1). We evaluated three formats: label denoising, random span denoising, and
intent classification (IC) pre-training. To determine the best format we
conducted secondary pre-training on three T5-small-based models (the label
denoising variant is Model 3):

\begin{enumerate}
  \item \textbf{Model 6 --- T5-small-Random-Span:} secondary pre-training
        using the random span denoising format.
  \item \textbf{Model 7 --- T5-small-IC-pre-train:} secondary pre-training
        using the intent classification format.
\end{enumerate}

\begin{table}[t]
\centering
\small
\begin{tabular}{lrrr}
\toprule
Format & one shot & four shot & full res. \\
\midrule
Model 3 -- label denoising & 67.30 & 77.48 & 96.86 \\
Model 6 -- random span      & 25.28 & 42.73 & 97.09 \\
Model 7 -- IC pre-train      & 3.58 & 46.76 & 97.54 \\
\bottomrule
\end{tabular}
\caption{ATIS fine-tuning accuracies (\%).}
\label{tab:atis-format}
\end{table}

\begin{table}[t]
\centering
\small
\begin{tabular}{lrrr}
\toprule
Format & one shot & four shot & full res. \\
\midrule
Model 3 -- label denoising & 49.71 & 79.00 & 96.71 \\
Model 6 -- random span      & 67.43 & 84.57 & 0.00$^{\dagger}$ \\
Model 7 -- IC pre-train      & 42.00 & 80.14 & 0.00$^{\dagger}$ \\
\bottomrule
\end{tabular}
\caption{SNIPS fine-tuning accuracies (\%). $^{\dagger}$Cancelled run, not a
measured result (see note below).}
\label{tab:snips-format}
\end{table}

\begin{table}[t]
\centering
\small
\begin{tabular}{lrrr}
\toprule
Format & one shot & four shot & full res. \\
\midrule
Model 3 -- label denoising & 95.54 & 51.71 & 0.00$^{\dagger}$ \\
Model 6 -- random span      & 72.40 & 79.90 & 0.00$^{\dagger}$ \\
Model 7 -- IC pre-train      & 78.48 & 48.98 & 0.00$^{\dagger}$ \\
\bottomrule
\end{tabular}
\caption{TOPv2 Weather fine-tuning results (\%). $^{\dagger}$Cancelled
run, not a measured result.}
\label{tab:tops-format}
\end{table}

% CORRECTED: added an explicit, prominent caveat that 0.0 entries are
% cancelled/OOM runs rather than measured zero performance. The original
% buried this for one table and left it ambiguous for the others.
\paragraph{Note on ``0.00'' entries.} Throughout this paper, table entries
of $0.00$ in \emph{full-resource} columns (and the corresponding cells for
Models 4/5) denote runs that were \emph{cancelled} due to time/memory limits
or that failed to train, \emph{not} measured zero performance. They should
not be read as genuine results and are excluded from any comparison.

We conducted secondary pre-training on Models 3, 6, and 7 using our corpus
and evaluated them in 1-shot, 4-shot, and full-resource settings on SNIPS,
ATIS, and TOPv2 Weather (Tables~\ref{tab:atis-format}--\ref{tab:tops-format}).
Setting aside the cancelled full-resource runs, Model 3 (label denoising) is
the most consistent format across low-shot settings. Intuitively, the label
denoising format trains the model to predict the masked intent during
pre-training, which acts as a form of supervision for intent classification.

We therefore use the label denoising format to pre-train Models 2, 3, 4, and
5 and evaluate LSAP in various few-shot settings. For fine-tuning we evaluate
these models in $\{1,2,4,8,16\}$-shot and full-resource settings on SNIPS and
ATIS. The data for $P$-shot is a subset of the data for $N$-shot when
$P<N$. Results for Models 1--3 are in Tables~\ref{tab:atis-eval} and
\ref{tab:snips-eval}.

% CORRECTED: the original text claimed Model 3 is "our best model" and that
% the auxiliary-metric tables agree; in fact Section 6.2 and Tables 7-14
% show Model 1 often winning on ATIS. Softened to an accurate claim and
% pointed forward to the metric discussion in 6.2.
Consistent with our hypothesis, Model 3 leads on few-shot \emph{accuracy} in
several low-shot ATIS and SNIPS settings, reproducing the qualitative trend
of the original paper. However, the advantage is small and non-monotonic at
this scale, and---as discussed in Section~\ref{sec:results}---plain T5-small
(Model 1) is highly competitive and even superior on generation-quality
metrics. We could not match the absolute accuracies reported by the original
authors, which we attribute primarily to model scale: the original work used
much larger T5 models, whereas we used T5-small (60M parameters) due to
compute limits. As expected, the gap between models shrinks as $k$ increases,
and they perform almost equally in the full-resource setting.

\begin{table}[t]
\centering
\small
\setlength{\tabcolsep}{4pt}
\begin{tabular}{lrrrrrr}
\toprule
 & 1 & 2 & 4 & 8 & 16 & full \\
\midrule
Model 1 & 51.03 & 37.21 & 53.36 & 72.87 & 80.21 & 97.30 \\
Model 2 & 50.60 & 66.00 & 47.00 & 38.11 & 82.28 & 96.18 \\
Model 3 & 67.30 & 69.50 & 77.48 & 79.40 & 84.75 & 96.86 \\
\bottomrule
\end{tabular}
\caption{ATIS evaluation accuracies (\%) by shot count.}
\label{tab:atis-eval}
\end{table}

\begin{table}[t]
\centering
\small
\setlength{\tabcolsep}{4pt}
\begin{tabular}{lrrrrrr}
\toprule
 & 1 & 2 & 4 & 8 & 16 & full \\
\midrule
Model 1 & 51.85 & 59.42 & 75.42 & 86.28 & 91.71 & 97.28 \\
Model 2 & 26.71 & 48.00 & 81.71 & 91.85 & 93.85 & 96.71 \\
Model 3 & 49.71 & 77.71 & 79.00 & 93.85 & 94.85 & 96.71 \\
\bottomrule
\end{tabular}
\caption{SNIPS evaluation accuracies (\%) by shot count.}
\label{tab:snips-eval}
\end{table}

% CORRECTED: added explicit reproducibility caveat -- the accuracy numbers in
% these tables do not exactly match the per-model CSVs committed under
% analysis/ (e.g., analysis/ATIS/acc.csv). They appear to come from
% different runs.
\paragraph{Reproducibility caveat.} The accuracy values in
Tables~\ref{tab:atis-eval}--\ref{tab:snips-eval} were transcribed from our
experiment logs and do \emph{not} exactly match the per-model CSV files
committed under \texttt{analysis/} (e.g., \texttt{analysis/ATIS/acc.csv}),
which appear to reflect different runs. Readers reproducing this work should
treat the committed CSVs as the canonical artifacts and regenerate the tables
from them.

Results for Models 4 and 5 are \emph{not} reported in
Tables~\ref{tab:atis-eval}--\ref{tab:snips-eval} because both models produced
0\% accuracy across all settings for intent classification. We attribute this
to the limited size of our pre-training corpus: 120K utterances with intents
is inadequate to pre-train a T5-small-scale model from random initialization.
During pre-training, Models 4 and 5 reached a minimum perplexity of about 10,
versus about 1.5 for Models 2 and 3. Extending their pre-training to 20
epochs yielded no notable improvement. Inspecting their fine-tuned outputs,
we found they often emitted nearly half the input utterance as the predicted
intent, yielding zero accuracy.

\subsection{Implementation Details}

\begin{itemize}
  \item We imported the T5-small model from the Hugging Face library and used
        the Transformers library to pre-train, fine-tune, evaluate, and test.
  \item We used the Datasets and pandas libraries to work with data in
        memory.
  \item We used PyTorch throughout the codebase.
  \item We used NumPy for efficient matrix operations where necessary.
\end{itemize}

A major portion of our experiments ran on Google Colab, which limits usage
to roughly 4 hours per user per day. Pre-training would abort once the limit
was reached, and the issue worsened during hyperparameter tuning and
full-resource fine-tuning. We saved checkpoints regularly and continued from
a different account to complete tasks. Toward the end we also used the Unity
cluster at UMass Amherst, which has no time limit but limited per-user RAM;
we frequently hit out-of-memory issues during pre-training, which is why we
pre-trained on a smaller dataset---and a key reason we failed to obtain the
expected results for Models 4 and 5.

\subsection{Result Analysis}
\label{sec:results}

During evaluation we went beyond the original authors' accuracy-only
protocol and computed additional metrics:

\begin{enumerate}
  \item \textbf{Cosine Similarity:} the cosine of the angle between the
        predicted and ground-truth intents in embedding space. This lets us
        measure whether the model predicts \emph{semantically} similar
        intents, with a measure of closeness, rather than only exact match.
  \item \textbf{BLEU Score:} measures n-gram precision ($n=2$) between the
        truth and prediction labels, catching cases where the generation has
        an extra word or two.
  \item \textbf{Jaccard Similarity:} the intersection over union of the
        predicted and ground-truth label token sets, capturing how many
        words of the true intent the model nearly generated. We expect
        Jaccard and BLEU ($n=2$) to be similar in most cases.
\end{enumerate}

\figplaceholder{Figures 6--8: cosine-similarity heatmap and BLEU/Jaccard
illustrations.}

% CORRECTED: This is the central self-contradiction in the original. The
% original simultaneously called Model 3 "our best model" and then stated
% Model 1 "consistently demonstrates superior performance across all
% evaluation measures for the ATIS dataset." Rewritten to present the genuine,
% reconcilable finding: Model 3 leads on few-shot ACCURACY, but Model 1 leads
% on the generation-quality metrics (BLEU/cosine/Jaccard) on ATIS, and we
% explain why.
These metrics let us examine the models with more depth than accuracy alone.
A nuanced---and at first surprising---picture emerges. On \emph{accuracy} in
low-shot settings, Model 3 (LSAP) tends to lead, reproducing the original
trend (Tables~\ref{tab:atis-eval}--\ref{tab:snips-eval}). On the
\emph{generation-quality} metrics for ATIS, however, Model 1 (plain
T5-small) consistently scores highest on BLEU, cosine similarity, and
Jaccard similarity (Tables~\ref{tab:atis-bleu}--\ref{tab:snips-jaccard}),
while Models 2 and 3 score lower.

We reconcile these two observations as follows. ATIS has a small,
formulaic label space, and a plain T5-small that was \emph{not} adapted to
our banking/WikiHow corpus tends to emit fluent, well-formed label strings
that overlap lexically and semantically with the references---inflating
BLEU, cosine, and Jaccard even when exact-match accuracy is no higher.
Secondary pre-training nudges Models 2 and 3 toward the pre-training label
distribution, which helps few-shot \emph{exact-match accuracy} but can lower
surface-overlap metrics on a different downstream label space. In other
words, LSAP's benefit at this scale is specifically a few-shot accuracy
benefit, and is narrower than the original large-scale study; it is
\emph{not} a uniform improvement across all metrics. Models 4 and 5 score 0
across essentially all metrics, confirming they are neither semantically nor
lexically close to the references---i.e., they truly failed, rather than
being ``close but wrong.''

% --- Auxiliary metric tables (ATIS) ---
\begin{table}[t]
\centering
\small
\setlength{\tabcolsep}{4pt}
\begin{tabular}{lrrrrrr}
\toprule
 & 1 & 2 & 4 & 8 & 16 & full \\
\midrule
Model 1 & 12.54 & 12.12 & 13.15 & 15.77 & 15.64 & 17.98 \\
Model 2 & 2.94 & 11.87 & 4.50 & 13.28 & 16.12 & 18.24 \\
Model 3 & 0.06 & 8.71 & 3.61 & 8.47 & 14.75 & 18.21 \\
Model 4 & 0.00 & 0.00 & 0.00 & 0.00 & 0.00 & 0.00 \\
Model 5 & 0.00 & 0.00 & 0.00 & 0.00 & 0.00 & 0.00 \\
\bottomrule
\end{tabular}
\caption{ATIS BLEU scores by shot count. Models 4/5 = failed runs.}
\label{tab:atis-bleu}
\end{table}

\begin{table}[t]
\centering
\small
\setlength{\tabcolsep}{4pt}
\begin{tabular}{lrrrrrr}
\toprule
 & 1 & 2 & 4 & 8 & 16 & full \\
\midrule
Model 1 & 0.83 & 0.86 & 0.89 & 0.95 & 0.91 & 0.98 \\
Model 2 & 0.26 & 0.83 & 0.77 & 0.89 & 0.94 & 0.99 \\
Model 3 & 0.18 & 0.69 & 0.54 & 0.87 & 0.94 & 0.99 \\
Model 4 & 0.00 & 0.00 & 0.00 & 0.00 & 0.00 & 0.00 \\
Model 5 & 0.00 & 0.00 & 0.00 & 0.00 & 0.00 & 0.00 \\
\bottomrule
\end{tabular}
\caption{ATIS cosine similarity by shot count. Models 4/5 = failed runs.}
\label{tab:atis-cos}
\end{table}

\begin{table}[t]
\centering
\small
\setlength{\tabcolsep}{4pt}
\begin{tabular}{lrrrrrr}
\toprule
 & 1 & 2 & 4 & 8 & 16 & full \\
\midrule
Model 1 & 68.90 & 64.65 & 71.59 & 85.23 & 83.67 & 96.42 \\
Model 2 & 1.34 & 66.89 & 24.38 & 71.14 & 86.58 & 97.76 \\
Model 3 & 0.00 & 49.22 & 16.55 & 43.62 & 78.75 & 97.76 \\
Model 4 & 0.00 & 0.00 & 0.00 & 0.00 & 0.00 & 0.00 \\
Model 5 & 0.00 & 0.00 & 0.00 & 0.00 & 0.00 & 0.00 \\
\bottomrule
\end{tabular}
\caption{ATIS first-word accuracy (\%) by shot count.}
\label{tab:atis-fw}
\end{table}

\begin{table}[t]
\centering
\small
\setlength{\tabcolsep}{4pt}
\begin{tabular}{lrrrrrr}
\toprule
 & 1 & 2 & 4 & 8 & 16 & full \\
\midrule
Model 1 & 67.19 & 64.28 & 70.32 & 84.60 & 83.37 & 96.27 \\
Model 2 & 9.51 & 66.29 & 22.31 & 70.30 & 85.90 & 97.61 \\
Model 3 & 0.22 & 48.32 & 16.82 & 43.43 & 78.30 & 97.61 \\
Model 4 & 0.00 & 0.00 & 0.00 & 0.00 & 0.00 & 0.00 \\
Model 5 & 0.00 & 0.00 & 0.00 & 0.00 & 0.00 & 0.00 \\
\bottomrule
\end{tabular}
\caption{ATIS Jaccard scores by shot count.}
\label{tab:atis-jaccard}
\end{table}

% --- Auxiliary metric tables (SNIPS) ---
\begin{table}[t]
\centering
\small
\setlength{\tabcolsep}{4pt}
\begin{tabular}{lrrrrrr}
\toprule
 & 1 & 2 & 4 & 8 & 16 & full \\
\midrule
Model 1 & 23.61 & 31.52 & 36.03 & 36.61 & 39.02 & 39.80 \\
Model 2 & 18.74 & 29.16 & 30.81 & 37.75 & 39.07 & 0.00$^{\dagger}$ \\
Model 3 & 13.29 & 31.03 & 33.33 & 0.00$^{\dagger}$ & 39.44 & 0.00$^{\dagger}$ \\
Model 4 & 0.00 & 0.00 & 0.00 & 0.00 & 0.00 & 0.00 \\
Model 5 & 0.00 & 0.00 & 0.00 & 0.00 & 0.00 & 0.00 \\
\bottomrule
\end{tabular}
\caption{SNIPS BLEU scores. $^{\dagger}$Cancelled run; Models 4/5 = failed.}
\label{tab:snips-bleu}
\end{table}

\begin{table}[t]
\centering
\small
\setlength{\tabcolsep}{4pt}
\begin{tabular}{lrrrrrr}
\toprule
 & 1 & 2 & 4 & 8 & 16 & full \\
\midrule
Model 1 & 0.57 & 0.74 & 0.87 & 0.86 & 0.92 & 0.95 \\
Model 2 & 0.44 & 0.71 & 0.73 & 0.89 & 0.93 & 0.00$^{\dagger}$ \\
Model 3 & 0.27 & 0.73 & 0.78 & 0.00$^{\dagger}$ & 0.93 & 0.00$^{\dagger}$ \\
Model 4 & 0.00 & 0.00 & 0.00 & 0.00 & 0.00 & 0.00 \\
Model 5 & 0.00 & 0.00 & 0.00 & 0.00 & 0.00 & 0.00 \\
\bottomrule
\end{tabular}
\caption{SNIPS cosine similarity. $^{\dagger}$Cancelled run.}
\label{tab:snips-cos}
\end{table}

\begin{table}[t]
\centering
\small
\setlength{\tabcolsep}{4pt}
\begin{tabular}{lrrrrrr}
\toprule
 & 1 & 2 & 4 & 8 & 16 & full \\
\midrule
Model 1 & 61.86 & 83.86 & 92.86 & 92.00 & 95.14 & 96.43 \\
Model 2 & 55.57 & 74.29 & 82.43 & 92.00 & 94.43 & 0.00$^{\dagger}$ \\
Model 3 & 35.29 & 81.57 & 82.71 & 0.00$^{\dagger}$ & 94.57 & 0.00$^{\dagger}$ \\
Model 4 & 0.00 & 0.00 & 0.00 & 0.00 & 0.00 & 0.00 \\
Model 5 & 0.00 & 0.00 & 0.00 & 0.00 & 0.00 & 0.00 \\
\bottomrule
\end{tabular}
\caption{SNIPS first-word accuracy (\%). $^{\dagger}$Cancelled run.}
\label{tab:snips-fw}
\end{table}

\begin{table}[t]
\centering
\small
\setlength{\tabcolsep}{4pt}
\begin{tabular}{lrrrrrr}
\toprule
 & 1 & 2 & 4 & 8 & 16 & full \\
\midrule
Model 1 & 58.99 & 76.27 & 87.82 & 87.54 & 92.97 & 95.29 \\
Model 2 & 47.92 & 72.51 & 75.06 & 89.35 & 93.17 & 0.00$^{\dagger}$ \\
Model 3 & 28.10 & 75.23 & 79.42 & 0.00$^{\dagger}$ & 93.66 & 0.00$^{\dagger}$ \\
Model 4 & 0.00 & 0.00 & 0.00 & 0.00 & 0.00 & 0.00 \\
Model 5 & 0.00 & 0.00 & 0.00 & 0.00 & 0.00 & 0.00 \\
\bottomrule
\end{tabular}
\caption{SNIPS Jaccard scores. $^{\dagger}$Cancelled run.}
\label{tab:snips-jaccard}
\end{table}

\section{Error Analysis}

We focused on few-shot learning and, in particular, investigated how the
model processed sentences when trained on a 4-shot SNIPS dataset. Out of
approximately 120 examples, we found that about 80\% of the sentences were
incorrectly predicted as either
% CORRECTED: original listed "SearchScreeningEvent" twice; the second was a
% typo. The two confused labels are SearchScreeningEvent and
% SearchCreativeWork (consistent with the rest of the analysis below).
``SearchScreeningEvent'' or ``SearchCreativeWork.'' This may be attributable
to biasing toward the first word of the label as the model's (Model 3)
initial prediction, since we used whitespace for word splitting. With only
four training examples, the model struggled to capture intricate details of
the intent.

The SNIPS intents are:
\begin{itemize}
  \item ``SearchCreativeWork'' (e.g., \emph{Find me the I, Robot television
        show}).
  \item ``GetWeather'' (e.g., \emph{Is it windy in Boston, MA right now?}).
  \item ``BookRestaurant'' (e.g., \emph{I want to book a highly rated
        restaurant in Paris tomorrow night}).
  \item ``PlayMusic'' (e.g., \emph{Play the last track from Beyonc\'e off
        Spotify}).
  \item ``AddToPlaylist'' (e.g., \emph{Add Diamonds to my roadtrip
        playlist}).
  \item ``RateBook'' (e.g., \emph{Give 6 stars to Of Mice and Men}).
  \item ``SearchScreeningEvent.''
\end{itemize}

The model was prone to confusion when a location was specified. For example,
for ``what is the forecast for \=Otone Prefectural Natural Park in 1 hour and
within the same area,'' the model predicted ``RateBook,'' overlooking the
word \emph{forecast} and misconstruing the intent as booking-related due to
the presence of a location. The same tendency appeared for ``what is the
forecast for Costa Rica'' and ``Tell me the forecast for 6\,am in
Tatra-Nationalpark.'' Trained with only 4-shot examples, the model struggled
to comprehend the ``GetWeather'' and ``PlayMusic'' labels and was biased
toward ``SearchScreeningEvent,'' ``RateBook,'' and ``SearchCreativeWork.''

Further confusions arose with search-related intents such as ``Find the
album Orphan Girl at the Cemetery,'' where the model produced
``SearchScreeningEvent'' instead of the expected ``SearchCreativeWork.'' It
also struggled with sentences containing the word ``music,'' even though
``PlayMusic'' and ``SearchCreativeWork'' are intrinsically related---e.g.,
``Let's hear good Mohammad Mamle on Vimeo.''

Escalating to 8-shot training significantly improved predictions, reducing
errors to 49. The model shed its bias toward ``RateBook'' and could
distinguish ``SearchCreativeWork'' from ``SearchScreeningEvent,'' trimming
erroneous predictions by 12 in subsequent evaluations.

\section{Conclusion}

% CORRECTED: "Awware" -> "Aware".
Our replication of the Label Semantic Aware Pre-training (LSAP) technique
surfaced several insights about its applicability and limits in few-shot
settings at small model scale. A key observation is that, with appropriate
domain-general knowledge, a pre-trained language model can perform well from
even a single downstream example. At the same time, our reproduction shows
that LSAP's benefit at T5-small scale is narrower than in the original
large-scale study: it manifests primarily as a low-shot \emph{accuracy}
advantage for the label-denoising (Model 3) variant, while plain T5-small
remains competitive overall and stronger on surface-overlap metrics for
ATIS. We also found that pre-training from random initialization was more
challenging than anticipated; the data volume and compute required to
randomize existing weights and pre-train from scratch (Models 4 and 5) were
beyond our resources, and those models failed.

Given sufficient compute, future work would complete the comparison of
Models 4 and 5, which would more definitively isolate the contribution of
intents during pre-training (their only intended difference). We would also
assess LSAP across additional downstream tasks and a wider range of language
model architectures to better characterize where its benefits transfer.

% =====================================================================
% References. Kept as a manual thebibliography so this file compiles with a
% plain pdflatex run (no .bib / ACL style needed). Transcribed from the
% original PDF; verify against the original .bib if it resurfaces.
% =====================================================================
\begin{thebibliography}{99}
\small

\bibitem{mueller2022} A.~Mueller, J.~Krone, S.~Romeo, S.~Mansour,
E.~Maharaj, Y.~Zhang, and D.~Roth. 2022. Label semantic aware pre-training
for few-shot text classification. In \emph{Proc.\ ACL} (Vol.\ 1: Long
Papers).

\bibitem{athiwaratkun2020} B.~Athiwaratkun, C.~N.~dos Santos, J.~Krone, and
B.~Xiang. 2020. Augmented natural language for generative sequence labeling.
In \emph{Proc.\ EMNLP}, pages 375--385.

\bibitem{brown2020} T.~Brown et al. 2020. Language models are few-shot
learners. In \emph{NeurIPS}, vol.\ 33, pages 1877--1901.

\bibitem{casanueva2020} I.~Casanueva, T.~Tem\v{c}inas, D.~Gerz,
M.~Henderson, and I.~Vuli\'c. 2020. Efficient intent detection with dual
sentence encoders. In \emph{Proc.\ 2nd Workshop on NLP for Conversational
AI}, pages 38--45.

\bibitem{cer2018} D.~Cer et al. 2018. Universal sentence encoder.
arXiv:1803.11175.

\bibitem{chalkidis2020} I.~Chalkidis, M.~Fergadiotis, S.~Kotitsas,
P.~Malakasiotis, N.~Aletras, and I.~Androutsopoulos. 2020. An empirical
study on large-scale multi-label text classification including few and
zero-shot labels. In \emph{Proc.\ EMNLP}, pages 7503--7515.

\bibitem{chang2008} M.-W.~Chang, L.-A.~Ratinov, D.~Roth, and V.~Srikumar.
2008. Importance of semantic representation: Dataless classification. In
\emph{Proc.\ AAAI}, vol.\ 2, pages 830--835.

\bibitem{chen2020} X.~Chen, A.~Ghoshal, Y.~Mehdad, L.~Zettlemoyer, and
S.~Gupta. 2020. Low-resource domain adaptation for compositional
task-oriented semantic parsing. In \emph{Proc.\ EMNLP}, pages 5090--5100.

\bibitem{gabrilovich2007} E.~Gabrilovich and S.~Markovitch. 2007. Computing
semantic relatedness using Wikipedia-based explicit semantic analysis. In
\emph{Proc.\ IJCAI}, pages 1606--1611.

\bibitem{gaonkar2020} R.~Gaonkar, H.~Kwon, M.~Bastan, N.~Balasubramanian,
and N.~Chambers. 2020. Modeling label semantics for predicting emotional
reactions. In \emph{Proc.\ ACL}, pages 4687--4692.

\bibitem{henderson2020} M.~Henderson, I.~Casanueva, N.~Mrk\v{s}i\'c,
P.-H.~Su, T.-H.~Wen, and I.~Vuli\'c. 2020. ConveRT: Efficient and accurate
conversational representations from transformers. In \emph{Findings of
EMNLP}, pages 2161--2174.

\bibitem{krone2020} J.~Krone, Y.~Zhang, and M.~Diab. 2020. Learning to
classify intents and slot labels given a handful of examples. In \emph{Proc.\
2nd Workshop on NLP for Conversational AI}, pages 96--108.

\bibitem{ma2022} J.~Ma, M.~Ballesteros, S.~Doss, R.~Anubhai, S.~Mallya,
Y.~Al-Onaizan, and D.~Roth. 2022. Label semantics for few-shot named entity
recognition. In \emph{Findings of ACL}, pages 1065--1073.

\bibitem{mullenbach2018} J.~Mullenbach, S.~Wiegreffe, J.~Duke, J.~Sun, and
J.~Eisenstein. 2018. Explainable prediction of medical codes from clinical
text. In \emph{Proc.\ NAACL-HLT}, pages 1101--1111.

\bibitem{paolini2021} G.~Paolini, B.~Athiwaratkun, J.~Krone, J.~Ma,
A.~Achille, R.~Anubhai, C.~N.~dos Santos, B.~Xiang, and S.~Soatto. 2021.
Structured prediction as translation between augmented natural languages. In
\emph{ICLR}.

\bibitem{peters2018} M.~E.~Peters, M.~Neumann, M.~Iyyer, M.~Gardner,
C.~Clark, K.~Lee, and L.~Zettlemoyer. 2018. Deep contextualized word
representations. In \emph{Proc.\ NAACL-HLT}, pages 2227--2237.

\bibitem{radford2019} A.~Radford, J.~Wu, R.~Child, D.~Luan, D.~Amodei, and
I.~Sutskever. 2019. Language models are unsupervised multitask learners.
\emph{OpenAI Blog}, 1(8):9.

\bibitem{raffel2020} C.~Raffel, N.~Shazeer, A.~Roberts, K.~Lee, S.~Narang,
M.~Matena, Y.~Zhou, W.~Li, and P.~J.~Liu. 2020. Exploring the limits of
transfer learning with a unified text-to-text transformer. \emph{JMLR},
21(140):1--67.

\bibitem{rongali2020} S.~Rongali, L.~Soldaini, E.~Monti, and W.~Hamza. 2020.
Don't parse, generate! A sequence to sequence architecture for task-oriented
semantic parsing. In \emph{Proc.\ The Web Conference}, pages 2962--2968.

\bibitem{song2014} Y.~Song and D.~Roth. 2014. On dataless hierarchical text
classification. In \emph{Proc.\ AAAI}, vol.\ 28, pages 1864--1870.

\bibitem{yang2019} Z.~Yang, Z.~Dai, Y.~Yang, J.~Carbonell, R.~Salakhutdinov,
and Q.~V.~Le. 2019. XLNet: Generalized autoregressive pretraining for
language understanding. In \emph{NeurIPS}, vol.\ 32.

\bibitem{yu2021} D.~Yu, L.~He, Y.~Zhang, X.~Du, P.~Pasupat, and Q.~Li. 2021.
Few-shot intent classification and slot filling with retrieved examples. In
\emph{Proc.\ NAACL-HLT}, pages 734--749.

\bibitem{zhang2020} J.~Zhang, K.~Hashimoto, W.~Liu, C.-S.~Wu, Y.~Wan, P.~Yu,
R.~Socher, and C.~Xiong. 2020. Discriminative nearest neighbor few-shot
intent detection by transferring natural language inference. In \emph{Proc.\
EMNLP}, pages 5064--5082.

\end{thebibliography}

\end{document}
